\chapter{Python Simulation: MASFE MDP Scheduler}
\label{appendix:code}

This appendix contains the key components of the MASFE Python simulation submitted as Code Submission Component~3. The full source file (\texttt{masfe\_simulation.py}) is available in the linked GitHub repository. The simulation runs with only \texttt{numpy} as a dependency (\texttt{pip install numpy}).

\medskip
\noindent Running the simulation produces the three-way comparison output (MASFE MDP vs fixed-schedule onboard vs naive raw-dump baseline) and exports the JSON metrics block used in the paper. Verified output:

\begin{itemize}
  \item Downlink reduction vs raw: \textbf{97\,\%}
  \item Energy saving vs raw: \textbf{21\,\%}
  \item Science retention: \textbf{100\,\%}
  \item Peak power: \textbf{9.6\,W} $\leq$ 14\,W budget \quad PASS
  \item Minimum battery SOC: \textbf{86\,\%} $>$ 15\,\% floor \quad PASS
\end{itemize}

%% -----------------------------------------------
\section*{CSC Fusion Function}
%% -----------------------------------------------

\begin{lstlisting}[language=Python, caption={Onboard Crop Stress Composite (CSC) fusion from MOD13A1 EVI and MOD11A1 LST. Equations correspond to Eq.~\eqref{eq:csc} in the main paper.}]
def compute_csc(evi_t, lst_t, evi_base, lst_base):
    """
    Crop Stress Composite -- per-patch anomaly index.
    Inputs:
        evi_t     : observed EVI per patch (MOD13A1)
        lst_t     : observed LST per patch in K (MOD11A1)
        evi_base  : per-patch healthy-season EVI baseline
        lst_base  : per-patch healthy-season LST baseline
    Returns:
        csc : array in [0, 1], 0=healthy, 1=severe stress
    """
    sigma_e = 0.042   # EVI noise + phenology (dimensionless)
    sigma_l = 1.20    # LST noise (K)

    evi_drop = np.maximum(0.0, (evi_base - evi_t) / sigma_e) / 5.0
    lst_rise = np.maximum(0.0, (lst_t - lst_base) / sigma_l) / 4.0

    csc = 0.55 * np.clip(evi_drop, 0, 1) \
        + 0.45 * np.clip(lst_rise, 0, 1)
    return np.clip(csc, 0, 1)
\end{lstlisting}

%% -----------------------------------------------
\section*{MDP Policy}
%% -----------------------------------------------

\begin{lstlisting}[language=Python, caption={MASFE two-stage MDP scheduling policy. Selects action from \{SKIP, MOD13, FUSE, FUSE\_PRIORITY\} based on battery SOC, EVI anomaly score, CSC level, and downlink window availability.}]
class MASFEPolicy:
    """
    Two-stage adaptive MDP scheduler.
    Stage 1 (MOD13): cheap EVI screening every pass.
    Stage 2 (FUSE): triggered by EVI anomaly or periodic
                    exploration (every explore_n steps).
    Degradation:
        SOC > 35% : full two-stage MDP (NOMINAL)
        15-35%    : Stage 1 (MOD13) only (CONSERVE)
        < 15%     : SKIP -- preserve battery (CRITICAL)
    """
    def __init__(self, evi_fuse_thr=0.18, csc_alert_thr=0.55,
                 explore_n=12, batt_crit=0.15, batt_cons=0.34):
        self.evi_fuse_thr  = evi_fuse_thr
        self.csc_alert_thr = csc_alert_thr
        self.explore_n     = explore_n
        self.batt_crit     = batt_crit
        self.batt_cons     = batt_cons

    def act(self, state):
        soc     = state.soc
        max_anom = state.evi_anom.max()
        max_csc  = state.csc.max()

        # CRITICAL: preserve battery
        if soc < self.batt_crit:
            return 'SKIP'

        # CONSERVE: Stage 1 screening only
        if soc < self.batt_cons:
            return 'MOD13'

        # Confirmed alert + downlink window -> priority TX
        if max_csc >= self.csc_alert_thr and state.downlink:
            return 'FUSE_PRIORITY'

        # EVI anomaly triggers Stage 2 confirmation
        if max_anom >= self.evi_fuse_thr:
            return 'FUSE'

        # High CSC but no downlink yet -> re-confirm, buffer
        if max_csc >= self.csc_alert_thr:
            return 'FUSE'

        # Periodic exploration (MDP discovery)
        if state.steps_since_fuse >= self.explore_n:
            return 'FUSE'

        # Default: cheap EVI screening
        return 'MOD13'
\end{lstlisting}

%% -----------------------------------------------
\section*{ISA Atmosphere Helper (TU Delft Template Reference)}
%% -----------------------------------------------

The ISA atmosphere model (shown in Appendix~A of the TU Delft template) was used to verify the 500\,km orbit altitude for solar flux and eclipse fraction calculations used in the power budget:

\begin{lstlisting}[language=Python, caption={ISA Calculator used for orbit power budget verification. Adapted from TU Delft report template Appendix A.}]
import math
g0, R = 9.80665, 287.0
layer1 = [0, 288.15, 101325.0]
alt = [0, 11000, 20000, 32000, 47000, 51000, 71000, 86000]
a   = [-.0065, 0, .0010, .0028, 0, -.0028, -.0020]

def atmosphere(h):
    for i in range(len(alt) - 1):
        if h >= alt[i]:
            layer0 = layer1[:]
            layer1[0] = min(h, alt[i + 1])
            if a[i] != 0:
                layer1[1] = layer0[1] + a[i] * (layer1[0] - layer0[0])
                layer1[2] = layer0[2] * (layer1[1] / layer0[1]) \
                            ** (-g0 / (a[i] * R))
            else:
                layer1[2] = layer0[2] * math.exp(
                    (-g0 / (R * layer1[1])) * (layer1[0] - layer0[0]))
    T, rho, P = layer1[1], layer1[2]/(R*layer1[1]), layer1[2]
    return [T, rho, P, math.sqrt(1.4 * R * T)]
\end{lstlisting}
